\documentclass[11pt]{article}
\usepackage[margin=1in]{geometry}
\usepackage{amsmath,amssymb,amsthm,mathtools}
\usepackage{enumitem}
\usepackage{microtype}
\usepackage[hidelinks]{hyperref}
\usepackage{booktabs}
\usepackage{array}
\usepackage{xcolor}
\usepackage{fancyhdr}
\usepackage{titlesec}
\usepackage{docmute}

\pagestyle{fancy}
\fancyhf{}
\lhead{AMLD Recursive What-If Imagination}
\rhead{Version 2.0 --- GBG+}
\cfoot{\thepage}

\newtheorem{definition}{Definition}[section]
\newtheorem{theorem}[definition]{Theorem}
\newtheorem{proposition}[definition]{Proposition}
\newtheorem{corollary}[definition]{Corollary}
\newtheorem{remark}[definition]{Remark}
\newtheorem{principle}[definition]{Design Principle}

\newcommand{\BANK}{\ensuremath{\mathsf{BANK}}}
\newcommand{\MQ}{\ensuremath{\mathsf{MAPQUEST}}}
\newcommand{\AMLD}{\ensuremath{\mathsf{AMLD}}}
\newcommand{\V}{\ensuremath{\mathsf{V}}}
\newcommand{\W}{\mathcal{W}}
\newcommand{\Th}{\operatorname{Th}}
\newcommand{\Pred}{\operatorname{Pred}}
\newcommand{\Enabled}{\operatorname{Enabled}}
\newcommand{\Score}{\operatorname{Score}}
\newcommand{\Novelty}{\operatorname{Novelty}}
\newcommand{\Gain}{\operatorname{Gain}}
\newcommand{\Cost}{\operatorname{Cost}}
\newcommand{\Res}{\operatorname{Res}}

\title{\textbf{Recursive ``What If?'' Imagination in AMLD}\\[0.4em]
\large Chess-Style Counterfactual Search Guided by MAPQUEST\\[0.4em]
\normalsize Version 2.0 --- GBG+ Reciprocal Extension}
\author{Brian Tenneson}
\date{October 2026}

\begin{document}

% Version 2.0 deliberately preserves Version 1.0 in its original order.
% docmute suppresses the old preamble and document wrapper while retaining
% the body.  The new title metadata above is therefore used by the original
% \maketitle command.
\documentclass[11pt]{article}
\usepackage[margin=1in]{geometry}
\usepackage{amsmath,amssymb,amsthm,mathtools}
\usepackage{enumitem}
\usepackage{microtype}
\usepackage[hidelinks]{hyperref}
\usepackage{booktabs}
\usepackage{array}
\usepackage{xcolor}
\usepackage{fancyhdr}
\usepackage{titlesec}

\pagestyle{fancy}
\fancyhf{}
\lhead{AMLD Recursive What-If Imagination}
\rhead{MAPQUEST Navigation}
\cfoot{\thepage}

\newtheorem{definition}{Definition}[section]
\newtheorem{theorem}[definition]{Theorem}
\newtheorem{proposition}[definition]{Proposition}
\newtheorem{corollary}[definition]{Corollary}
\newtheorem{remark}[definition]{Remark}
\newtheorem{principle}[definition]{Design Principle}

\newcommand{\BANK}{\ensuremath{\mathsf{BANK}}}
\newcommand{\MQ}{\ensuremath{\mathsf{MAPQUEST}}}
\newcommand{\AMLD}{\ensuremath{\mathsf{AMLD}}}
\newcommand{\V}{\ensuremath{\mathsf{V}}}
\newcommand{\W}{\mathcal{W}}
\newcommand{\Th}{\operatorname{Th}}
\newcommand{\Pred}{\operatorname{Pred}}
\newcommand{\Enabled}{\operatorname{Enabled}}
\newcommand{\Score}{\operatorname{Score}}
\newcommand{\Novelty}{\operatorname{Novelty}}
\newcommand{\Gain}{\operatorname{Gain}}
\newcommand{\Cost}{\operatorname{Cost}}
\newcommand{\Res}{\operatorname{Res}}

\title{\textbf{Recursive ``What If?'' Imagination in AMLD}\\[0.4em]
\large Chess-Style Counterfactual Search Guided by MAPQUEST}
\author{Brian Tenneson}
\date{September 2026}

\begin{document}
\maketitle

\begin{abstract}
This note develops a precise notion of imagination for an AMLD-style multi-agent theorem-settling architecture. The motivating analogy is chess: a strong player or engine is imaginative not because it violates the rules, but because it recursively asks ``what if?'' about unusual continuations, hypothetical positions, sacrifices, and long-range plans. We transfer that idea to automated proof search. A certified proof state is stored in a monotone verifier-gated \BANK, while imaginative agents create speculative counterfactual states that never acquire theorem status merely by being imagined. MAPQUEST supplies the address system, inference geometry, target compass, and navigational machinery for exploring each speculative branch. The verifier remains sovereign: only verified consequences may be promoted from imagination into the certified mathematical state. This separation yields a rigorous form of machine imagination that can be arbitrarily adventurous in search while remaining conservative in proof.
\end{abstract}

\section{MAPQUEST as Navigation Toward a Winning Strategy}

A useful informal description is:
\begin{quote}
\textbf{MAPQUEST is the navigation system for finding a winning strategy.}
\end{quote}
In theorem proving, the ``winning strategy'' is a finite verifier-accepted route from the initial environment to a target well-formed formula (WFF) $T$.

The current MAPQUEST convention keeps each WFF intrinsically zero-dimensional. Formula addresses identify locations, while inference applications carry the geometric dimension induced by rule arity. If rule $r$ has arity $k$, its navigation component has the form
\[
X_r=\{r\}\times (\mathbb{Q}\cap[0,1])^k.
\]
Across heterogeneous rules, the global inference-action space is the tagged coproduct
\[
\bigsqcup_{r\in z}
\left(\{r\}\times(\mathbb{Q}\cap[0,1])^{\operatorname{ar}(r)}\right).
\]
Thus MAPQUEST is more than a move selector. It provides an address system, a geometry of legal inference applications, a multiscale traversal, a target-oriented compass, a memory of verified progress, and a route-to-certificate mechanism.

In chess language, the analogy is therefore not merely
\[
\text{MAPQUEST}=\text{choose a move}.
\]
A better analogy is
\begin{center}
\fbox{\parbox{0.88\textwidth}{\centering
\textbf{MAPQUEST is the navigable geometry and compass of move applications leading toward a winning line.}}}
\end{center}

\section{Imagination as Recursive ``What If?''}

The central proposal of this note is that imagination in AMLD should be defined recursively.

\begin{definition}[Recursive what-if imagination]
Let $B_t$ be the certified proof state at time $t$, let $T$ be the target, and let $d\in\mathbb{N}$ be a counterfactual depth. The depth-$d$ imaginative operator
\[
\W_d(B_t,T)
\]
constructs and explores hypothetical extensions of $B_t$ by repeatedly asking questions of the form
\[
\text{``What if candidate formula or inference }H\text{ were available?''}
\]
and, for each such hypothetical state, recursively asking further ``what if?'' questions up to depth $d$.
\end{definition}

At depth zero, no speculation is performed:
\[
\W_0(B,T)=B.
\]
At positive depth, candidate hypotheses $H_1,\dots,H_m$ induce speculative states
\[
\widetilde B^{(i)}=B\cup\{H_i\},
\]
and recursion continues from each such state:
\[
\W_d(B,T)
=
\bigcup_{i=1}^{m}
\W_{d-1}(\widetilde B^{(i)},T),
\]
subject to a finite search budget, admissibility filters, and MAPQUEST navigation priorities.

The formulas $H_i$ need not yet be theorems. They are hypotheses on an analysis board. Their role is to reveal useful future structure, not to alter the certified mathematical world.

\subsection{The chess interpretation}

A chess engine recursively explores
\[
S\xrightarrow{m_1}S_1\xrightarrow{m_2}S_2\xrightarrow{m_3}\cdots
\]
and repeatedly asks:
\begin{quote}
What if I play this move? What if the opponent replies that way? What if I sacrifice material here? What if I first create this position and only then attack?
\end{quote}
Imagination lies in the recursion and in the choice of which counterfactual branches deserve attention. The rules of chess remain fixed.

The AMLD analogue is
\[
B_t
\rightsquigarrow
\widetilde B_1
\rightsquigarrow
\widetilde B_2
\rightsquigarrow
\cdots,
\]
where $\rightsquigarrow$ denotes a speculative transition rather than a certified BANK deposit.

\section{Certified Reality Versus the Analysis Board}

The architecture must distinguish two kinds of state.

\paragraph{Certified state.}
\[
B_t\subseteq \Th_{\mathcal F}(\Gamma)
\]
contains only axioms, hypotheses, and verifier-accepted consequences.

\paragraph{Speculative state.}
\[
\widetilde B=B_t\cup H
\]
may contain unproved formulas introduced solely for counterfactual analysis.

The two must never be confused.

\begin{principle}[Verifier sovereignty]
No speculative formula becomes a certified BANK entry merely because it creates a promising route. Promotion from speculation to BANK requires an independently verifier-accepted derivation from the frozen environment.
\end{principle}

This gives the basic information flow
\[
\boxed{
\text{Imagine}
\longrightarrow
\text{Navigate}
\longrightarrow
\text{Select subgoal}
\longrightarrow
\text{Prove}
\longrightarrow
\text{Verify}
\longrightarrow
\BANK.
}
\]

\begin{theorem}[Counterfactual conservativity]
Assume the certified BANK initially contains only formulas derivable from $\Gamma$, every BANK deposit is accepted by a sound verifier, and speculative states cannot write directly to BANK. Then arbitrary finite or infinite counterfactual exploration cannot make BANK unsound.
\end{theorem}

\begin{proof}
Counterfactual search may generate any number of speculative formulas, but none enters BANK through speculation alone. By hypothesis, every actual BANK deposit is verifier accepted, and verifier soundness implies that every accepted deposit is derivable from $\Gamma$. Induction over BANK deposits therefore gives
\[
B_t\subseteq\Th_{\mathcal F}(\Gamma)
\]
for every certified stage $t$, independently of the behavior of the imaginative search layer.
\end{proof}

This theorem is the formal reason AMLD can be imaginative without weakening logical rigor.

\section{What MAPQUEST Does Inside Imagination}

MAPQUEST should not merely navigate the certified state. It should also navigate each speculative board.

Given a hypothetical state
\[
\widetilde B=B_t\cup\{H_1,\dots,H_j\},
\]
MAPQUEST can recompute or estimate:
\begin{enumerate}[label=(\roman*)]
\item backward target cones for $T$;
\item missing-premise counts for target-producing frames;
\item residual proof estimates such as $\widehat h(T\mid\widetilde B)$;
\item newly enabled inference neighborhoods;
\item Hilbert-address regions that become promising;
\item newly valuable intermediate WFFs;
\item routes that were invisible or unattractive from $B_t$ alone.
\end{enumerate}

Thus imagination is not detached fantasy. It is a controlled perturbation of the state followed by disciplined MAPQUEST navigation.

\section{High-Leverage Imagined Lemmas}

Suppose $d(T\mid B)$ is an estimated remaining proof distance. For an imagined lemma $L$, define its counterfactual leverage by
\[
\Delta_T(L\mid B_t)
=
 d(T\mid B_t)-d(T\mid B_t\cup\{L\}).
\]
A large positive value means that, \emph{if} $L$ were available, the target would appear substantially easier.

This does not prove $L$. It changes what AMLD should try to prove next.

\begin{definition}[High-leverage imagined lemma]
An imagined lemma $L$ is high leverage at time $t$ if
\[
\Delta_T(L\mid B_t)\ge \theta
\]
for a chosen threshold $\theta>0$, or if it produces a comparably large increase in target-relevant enabled inference structure.
\end{definition}

The corresponding strategy is
\[
\boxed{
\text{Imagine }L
\;\to\;
\text{measure its effect on navigation}
\;\to\;
\text{promote }L\text{ to a real subtarget}
\;\to\;
\text{attempt to prove }L.
}
\]

This is the proof-search analogue of a chess player imagining a knight on $f5$, discovering that the resulting position is winning, and then asking how to maneuver the knight there legally.

\section{Option Gain and the Mathematics of Creative Potential}

A move or lemma can be valuable because it creates many future possibilities, even when it does not immediately approach the target.

Define the enabled-inference set at state $B$ by
\[
\Enabled(B)=\{c:\text{all certified premises required by candidate }c\text{ are available in }B\}.
\]
For a hypothetical lemma $L$, define its option gain by
\[
\Gain(L\mid B_t)
=
\left|
\Enabled(B_t\cup\{L\})\setminus\Enabled(B_t)
\right|.
\]
A high option-gain lemma is analogous to a chess move that activates several pieces simultaneously.

A generic imagination score can therefore combine target leverage, novelty, option gain, and cost:
\[
\Score(L)
=
\alpha\,\Delta_T(L\mid B_t)
+
\beta\,\Novelty(L)
+
\gamma\,\Gain(L\mid B_t)
-
\delta\,\Cost(L).
\]
The coefficients may be static, learned, or adaptively controlled by AMLD's awareness layer.

\section{Strategic Sacrifice}

Chess imagination often includes sacrifice: accepting an immediate local disadvantage in exchange for a superior future position. AMLD should permit the analogous behavior.

Suppose candidate branch $c_1$ appears to reduce the current residual estimate immediately, while $c_2$ temporarily makes it worse but opens a substantially richer future inference neighborhood. A purely greedy controller selects $c_1$. An imaginative controller may select $c_2$.

Define a schematic sacrifice score
\[
\operatorname{Sac}(c)
=
\operatorname{FutureGain}(c)
-
\lambda\operatorname{ImmediateLoss}(c).
\]
Branches with positive sacrifice score are candidates for deliberate non-greedy exploration.

This is not permission to accept invalid mathematics. The ``loss'' is only a search or heuristic loss. The verifier remains unchanged.

\section{Recursive Depth and Counterfactual Trees}

The strongest form of imagination is not a single hypothetical lemma but nested counterfactual reasoning:
\[
\text{What if }L_1?
\]
then
\[
\text{If }L_1,\text{ what if }L_2?
\]
then
\[
\text{If }L_1,L_2,\text{ what if }L_3?
\]
and so on.

This creates a counterfactual tree or DAG. A node is a speculative state
\[
\widetilde B_{\sigma}
=
B_t\cup\{L_{\sigma_1},\dots,L_{\sigma_j}\},
\]
where $\sigma$ records the hypothetical choices along the branch.

A branch receives attention according to a priority such as
\[
Q(\sigma)
=
-\alpha\widehat h(T\mid\widetilde B_\sigma)
+\beta\operatorname{Novelty}(\sigma)
+\gamma\operatorname{OptionGain}(\sigma)
-\delta\operatorname{Cost}(\sigma).
\]

The imaginative process is therefore recursive but bounded:
\begin{enumerate}[label=(\arabic*)]
\item propose unusual but structurally admissible hypotheses;
\item extend the speculative state;
\item run MAPQUEST locally on that state;
\item evaluate target residual and future option gain;
\item recurse on the best or most novel branches;
\item extract high-leverage hypothetical lemmas;
\item return them to the certified system as subgoals, not as facts.
\end{enumerate}

\section{AMLD Agent Roles Under Recursive Imagination}

A multi-agent architecture can assign different counterfactual personalities:

\begin{center}
\begin{tabular}{>{\centering\arraybackslash}p{0.08\textwidth} p{0.78\textwidth}}
\toprule
Agent & Counterfactual role \\
\midrule
$P$ & asks what hypothetical lemmas or decompositions would make $T$ provable quickly; \\
$R$ & asks what hypothetical routes would make $\neg T$ provable or expose failure in the proposed proof direction; \\
$I$ & explores alternative assumption sets, model-like obstructions, or certificate structures relevant to independence protocols; \\
$C$ & challenges speculative branches for hidden inconsistency, provenance conflict, malformed transitions, or environment drift; \\
$D$ & Dreamer: deliberately proposes unusual, low-prior, high-novelty counterfactual branches; \\
$A$ & Analyst: measures how each imagined state changes residual distance, option gain, and target geometry. \\
\bottomrule
\end{tabular}
\end{center}

Each agent can run its own MAPQUEST navigator while sharing the same frozen formal environment and the same verifier-gated BANK.

\section{The Dreamer Depth Charge}

Imagination need not run at full intensity continuously. A natural policy is to trigger deep counterfactual recursion when ordinary search stagnates.

For example, if a residual statistic changes negligibly over a window,
\[
\widehat h_{t-k}(T)-\widehat h_t(T)\approx 0,
\]
AMLD may launch a \emph{Dreamer Depth Charge}:
\begin{enumerate}[label=(\alph*)]
\item checkpoint the certified state;
\item increase counterfactual depth;
\item invert or flatten ordinary move priorities;
\item explore low-frequency Hilbert regions;
\item hypothesize intermediate lemmas outside the current target cone;
\item evaluate sacrifices and option gain;
\item return only verified discoveries to BANK;
\item resume ordinary search with any new certified information.
\end{enumerate}

This formalizes a chess-like change of plan without changing the logical rules.

\section{A Minimal Recursive Algorithm}

A compact version of the imaginative controller is:

\begin{verbatim}
IMAGINE(B, T, depth, budget):
    if depth = 0 or budget exhausted:
        return []

    H := propose_counterfactuals(B, T)
    rank H using target leverage, novelty, option gain, and cost

    discoveries := []
    for L in selected H:
        B_tilde := B union {L}          # speculative only
        analysis := MAPQUEST(B_tilde, T, speculative=True)

        if analysis says L is high-leverage:
            submit L as a real subtarget to the certified prover

        discoveries += IMAGINE(B_tilde, T, depth-1, remaining_budget)

    return discoveries
\end{verbatim}

The critical line is the comment ``speculative only.'' The hypothetical extension is never a BANK deposit.

\section{Winning Strategy Versus Winning Line}

The chess analogy distinguishes a strategy from a certified line. MAPQUEST navigates toward a winning strategy, but it does not assume that the winning strategy is known in advance.

In theorem-proving notation,
\[
\Gamma
\longrightarrow
\text{MAPQUEST navigation}
\longrightarrow
\text{candidate proof strategy}
\longrightarrow
\Pi_T
\longrightarrow
\Gamma\vdash T.
\]

Recursive imagination enriches the middle term:
\[
\Gamma
\longrightarrow
\underbrace{\text{recursive ``what if?'' worlds}}_{\text{imagination}}
\longrightarrow
\underbrace{\text{MAPQUEST navigation in those worlds}}_{\text{geometry + compass}}
\longrightarrow
\text{valuable real subgoals}
\longrightarrow
\text{verified winning proof}.
\]

Thus the most concise conceptual statement is
\begin{center}
\fbox{\parbox{0.88\textwidth}{\centering
\textbf{Imagination recursively asks ``what if?''; MAPQUEST navigates the resulting possibilities; the verifier decides what is real.}}}
\end{center}

\section{Conclusion}

A good AMLD should not equate imagination with randomness, temperature, or unrestricted generation. Chess suggests a stronger definition: imagination is the recursive construction and evaluation of counterfactual continuations. In mathematics, this means temporarily supposing that strategically chosen lemmas, decompositions, or intermediate structures are available, then using MAPQUEST to determine how the geometry of possible proof applications changes.

The architecture remains logically conservative because speculative states are separated from certified state. The Dreamer may recursively ask ``what if?'' without limit in principle; the only route back to mathematical reality is through proof and verification.

The resulting design can be summarized as
\begin{center}
\fbox{\parbox{0.88\textwidth}{
\textbf{AMLD imagination} = recursive counterfactual construction
+ MAPQUEST navigation
+ target-leverage estimation
+ option-gain valuation
+ strategic sacrifice
+ promotion of valuable imagined lemmas to real subgoals
+ verifier-gated return to BANK.}}
\end{center}

This makes imagination operational: the system does not merely search harder. It repeatedly asks what world would make the target easier, studies that world, and then attempts to build the verified mathematical road that actually gets there.

\end{document}


\appendix
\section{GBG+ Realization and the Reciprocal Effect of Imagination}

This appendix adds a two-way connection between recursive AMLD imagination and the formal GBG/GBG+ substrate.  The first direction is conservative: GBG+ supplies a rule-complete legal-transition semantics for the machinery that creates, schedules, evaluates, and retires speculative proof worlds.  The second direction is more important conceptually: recursive imagination supplies a new optional capability for GBG+ itself.  A board-player need not merely schedule, remember, certify, hand off, or mutate a single current board.  It may legally create verifier-isolated counterfactual subgames, explore them recursively, extract strategic information, and return that information to the parent game without treating imagined states as certified states.

\subsection{GBG+ constrains imagination}

Let a rule-complete enlarged state contain
\[
S=(q,M_1,\ldots,M_N,M_B,\sigma,\Pi,\mathcal I),
\]
where $q$ is the current certified problem state, $M_i$ are local memories, $M_B$ is trusted board memory, $\sigma$ is scheduler state, $\Pi$ is provenance/certificate data, and $\mathcal I$ is a finite registry of speculative worlds.  Ordinary AMLD processes are represented by ordinary participants $P_1,\ldots,P_N$, while the distinguished board-player is
\[
P_0=B^\star.
\]
Every operation attributed to imagination---spawning a speculative world, extending it, scheduling an analysis process, checkpointing it, extracting a candidate lemma, retiring the world, or requesting verification---must be a declared legal transition of the enlarged game.  Thus recursive imagination inherits the GBG+ non-sovereignty principle: imagination may be adventurous about hypotheses, but the machine executing imagination has no undeclared meta-authority.

\begin{theorem}[GBG+ realization of recursive imagination]
Suppose an AMLD recursive-imagination architecture has finitely encoded rule-complete states, declared operational transitions, verifier-gated trusted deposits, and no transition from speculative state to trusted state except through an accepted certificate.  Then the operational architecture admits a GBG+ presentation in which each finite imaginative run is a legal path in an enlarged state graph.
\end{theorem}

\begin{proof}
Encode the complete operational state, including the speculative-world registry and all rule-relevant metadata, as the enlarged GBG state.  Give each declared AMLD operation a canonical move label and its uniquely determined partial successor rule.  Authorize ordinary search operations to the corresponding ordinary participants and board-level scheduling, routing, storage, checkpoint, spawn, retire, and verifier-dispatch operations to $P_0$ when their declared preconditions hold.  Rule completeness makes legal continuations state-determined, and the no-extra-authority condition ensures that the board-player acts only through declared rules.  Hence a finite run is exactly a finite legal transition history in the enlarged game.
\end{proof}

\subsection{Imagination enlarges what the plus can mean operationally}

The base GBG already generalizes ordinary board games by treating rule-complete states as formal objects and legal moves as declared partial transitions.  GBG+ adds the board as a distinguished participant.  Recursive imagination reveals a stronger use of that added participant: the board can operate not only \emph{on} the current game but also on a controlled family of hypothetical games derived from it.

This motivates an optional \emph{Imagination Package} for GBG+.

\begin{definition}[Counterfactual subgame]
Given a certified parent state $S$ and a finite hypothetical packet $H$, a counterfactual subgame is a tagged GBG+ instance
\[
G^+_{S,H}
\]
whose initial speculative state contains the parent snapshot together with $H$, and whose trusted-output channel is disabled except through an explicit verifier-gated return rule.
\end{definition}

The tag is essential.  It distinguishes
\[
\text{true in the certified parent game}
\qquad\text{from}\qquad
\text{assumed inside a counterfactual subgame}.
\]
The subgame may reason freely from its temporary assumptions.  Those assumptions acquire no certified status in the parent merely because the subgame finds them useful.

\begin{principle}[I-1: legal imaginative spawning]
A board-player may create a counterfactual subgame only through a declared spawn transition
\[
\operatorname{Spawn}(S,H):S\longrightarrow S',
\]
where $S'$ records the new tagged child and its resource budget.
\end{principle}

\begin{principle}[I-2: speculative isolation]
No proposition whose only support uses a hypothetical assumption in a child game may enter trusted parent memory as a theorem.
\end{principle}

\begin{principle}[I-3: verifier-gated return]
A child may return a candidate, heuristic score, proposed subgoal, route, obstruction, or search policy to the parent.  A mathematical proposition is promoted to trusted parent state only after the parent verifier accepts a certificate from the frozen certified environment.
\end{principle}

\begin{principle}[I-4: recursive spawning]
A counterfactual child may itself spawn counterfactual descendants when the selected GBG+ architecture declares such moves and charges their costs.  Thus imagination generates a finite or budget-bounded tree or DAG of GBG+ instances rather than an untyped collection of fantasies.
\end{principle}

\begin{principle}[I-5: provenance of hypothetical dependence]
Every returned object records whether it is certified, heuristic, or dependent on one or more hypothetical assumptions.  Erasing the speculative tag is not a legal handoff transformation.
\end{principle}

\subsection{The imaginative federation}

For a parent game $G^+$, let
\[
\mathfrak I_d(G^+,S,T)
\]
denote the family of counterfactual descendants explored to depth at most $d$ while pursuing target $T$.  Each node is itself a legal GBG+ instance or legal tagged local board.  Parent-child edges are spawn/handoff relations rather than proof implications.

This changes the geometry of search.  Ordinary GBG reachability asks which states can be reached by legal moves in one game.  Imaginative GBG+ additionally asks which \emph{legal game instances} can be generated, inspected, compared, and retired by the board-player.  The search object therefore has two levels:
\[
\text{within-game state graph}
\qquad\text{and}\qquad
\text{between-game imagination graph}.
\]
The second graph is generated by legal board actions.  It does not replace the first.

\begin{definition}[Imagination graph]
The imagination graph has a vertex for each active or archived tagged counterfactual game instance and a directed edge
\[
G^+_a\longrightarrow G^+_b
\]
when $G^+_b$ was legally spawned from $G^+_a$.  Edge metadata records the hypothetical packet, cost, provenance boundary, and return policy.
\end{definition}

This provides a formal home for the recursive ``what if?'' tree already used by AMLD.  It also gives GBG+ a natural recursive interpretation: a board-player may legally create boards that contain board-players that create further boards, subject to declared depth, budget, verifier, and handoff rules.

\subsection{Counterfactual conservativity across the federation}

\begin{theorem}[Federated counterfactual conservativity]
Assume every trusted parent deposit requires verifier acceptance from the certified parent environment and every hypothetical dependency remains tagged until independently discharged.  Then arbitrary finite recursive spawning of counterfactual GBG+ children does not enlarge the trusted parent consequence set merely by imagination.
\end{theorem}

\begin{proof}
Induct on trusted parent deposits rather than on the number of speculative descendants.  The base trusted state is certified.  A new trusted proposition can enter only through the verifier-gated return rule, so it has an accepted parent-environment certificate.  Child assumptions, regardless of recursive depth, cannot themselves satisfy this rule.  Therefore speculative branching can change search priorities and candidate selection but cannot by itself add an uncertified theorem to trusted parent state.
\end{proof}

The theorem is stronger operationally than a single-board statement because the number and depth of speculative boards are irrelevant to soundness.  Soundness is controlled at the return boundary.

\subsection{How imagination changes the interpretation of GBG+}

Recursive imagination suggests that the plus sign should be read as more than ``one extra player label.''  Formally, the minimal foundation may continue to define GBG+ by the distinguished board-player $P_0=B^\star$; that is the stable core.  Operationally, however, the board-player becomes the gateway to a family of conservative meta-capabilities:
\[
\boxed{
\begin{array}{c}
\text{schedule}\\
\text{remember}\\
\text{route and hand off}\\
\text{verify and certify}\\
\text{spawn hypothetical boards}\\
\text{compare counterfactual futures}\\
\text{recurse over boards}\\
\text{return only verifier-safe discoveries}
\end{array}}
\]
This is exactly where the generalized-board-game idea becomes especially powerful.  A ``board'' need not be a chessboard, a single theorem state, or even a single fixed search world.  It can be a rule-complete formal environment whose legal dynamics include the controlled construction of further formal environments.

The reciprocal relationship is therefore
\[
\boxed{
\text{GBG+ gives imagination legal semantics; imagination gives GBG+ recursive counterfactual reach.}
}
\]

\subsection{Interaction with MAPQUEST}

MAPQUEST can be lifted from navigation inside one state graph to navigation across the imagination graph.  Within a child game it estimates target geometry, residual distance, option gain, and useful local regions.  Across children it can rank whole counterfactual worlds by a score such as
\[
Q(G^+_a)=
-\alpha\widehat h(T\mid G^+_a)
+\beta\operatorname{Novelty}(G^+_a)
+\gamma\operatorname{OptionGain}(G^+_a)
-\delta\operatorname{Cost}(G^+_a).
\]
Thus MAPQUEST acquires two navigation scales: proof navigation inside a board and imagination navigation among boards.

\subsection{Version 2.0 summary}

Version 1.0 established the central separation
\[
\text{imagine}\to\text{navigate}\to\text{prove}\to\text{verify}\to\BANK.
\]
Version 2.0 places that pipeline inside GBG+ and then sends the influence back in the other direction.  The resulting architecture is
\[
\boxed{
\begin{array}{c}
\text{certified parent GBG+}\\
\downarrow\;\operatorname{Spawn}\\
\text{counterfactual GBG+ federation}\\
\downarrow\;\text{MAPQUEST analysis}\\
\text{candidate subgoals / strategies / obstructions}\\
\downarrow\;\text{independent proof and verification}\\
\text{trusted parent BANK}
\end{array}}
\]
The conceptual advance is that imagination is no longer merely a heuristic running inside AMLD.  It becomes a candidate optional GBG+ package: a verifier-safe mechanism by which a generalized board can legally create, explore, compare, and retire hypothetical generalized boards.

\end{document}
